\section{Introduction}

Reading is often described as “the process of gaining meaning from print” \citep[][p.~34]{raynerHowPsychologicalScience2001}. Although comprehension ultimately depends on sentence- and discourse-level integration, efficient access to individual word forms and meanings remains a foundational component of skilled reading. A central challenge for readers is therefore to extract structure from written words rapidly and flexibly, using information distributed across orthography, morphology, and semantics.

In languages such as English, this challenge is amplified by the properties of the writing system itself. English exhibits irregular and context-sensitive grapheme–phoneme mappings, rendering phonological decoding alone an insufficient basis for fluent word recognition. At the same time, English preserves regularities at higher representational levels, particularly in its morphological structure. Words that share a stem or affix often exhibit consistent orthographic patterns and systematic semantic relationships (e.g., heal–health, teach–teacher). As a result, skilled reading in English requires sensitivity not only to phonological correspondences but also to recurring morphological structure that links form and meaning.

The question of how morphological structure is represented and processed during word recognition has generated substantial empirical and theoretical debate. A core issue concerns the nature of the units involved: whether skilled readers decompose complex words into their constituent morphemes during recognition, access whole-word representations directly, or exploit learned statistical regularities that render such architectural distinctions unnecessary. These possibilities are not merely theoretical alternatives—they carry different implications for what it means to be a skilled morphological processor, and for how individual differences in reading ability should modulate the recognition of morphologically complex words. Adjudicating among them requires identifying variables that are selectively diagnostic of the proposed mechanisms, and two such variables—stem frequency and morphological family size—have proven particularly central to this debate.

\subsection{Stem Frequency and Morphological Family Size: Theoretical Interpretations Across Current Models}

Stem frequency refers to the summed token frequency of a base morpheme across all its surface realizations—inflected and derived forms alike—and is often operationalized as lemma frequency or cumulative root frequency [CITATION]. Morphological family size refers to the number of distinct word types that share a base morpheme, encompassing both compounds and derivatives [CITATION]. Crucially, the two variables are empirically dissociable: a base can be high in stem frequency but low in family size, or vice versa. Both facilitate visual lexical decision and naming, but their effects are independent when appropriately controlled [CITATION], and it is this dissociation that has driven theoretical debate. What each effect represents depends substantially on which processing architecture one adopts, and current models differ not only in their answers but in the kinds of representations they permit as answers.

In the classical framework developed by Baayen and Schreuder [CITATION] and related dual-route accounts, the two effects were treated as converging diagnostics of separate processing routes. Stem frequency was held to reflect the ease of accessing a stored morpheme entry via a decompositional route: because complex words are parsed into their constituent morphemes during recognition, the frequency of the stem as an independent lexical unit facilitates that parsing step and any subsequent form built upon it. Morphological family size, by contrast, was attributed to the whole-form route and to the richness of the semantic network surrounding a base morpheme: encountering a word activates not only its own representation but those of its morphological relatives, and a larger family produces broader co-activation of the base representation [CITATION]. On this view, the dissociation between the two effects was itself evidence for a dual-route architecture, with each variable selectively loading on one route.

Subsequent models have substantially revised these interpretations, though they differ in the depth of revision they undertake. The early automatic decomposition account developed by Rastle, Davis, and colleagues [CITATION], drawing on masked morphological priming, proposes that decomposition is not one route among several but an obligatory initial parsing step that applies to all morphologically structured—and apparently morphologically structured—inputs regardless of semantic transparency. On this account, stem frequency reflects the ease of this early parsing stage: higher stem frequency facilitates the activation of the morpheme representation that is the immediate output of decomposition. Morphological family size is then attributed to a later, semantically sensitive stage, consistent with the finding that opaque relatives contribute less to the family size effect than transparent ones [CITATION]. The two-stage sequential architecture thus preserves the dissociation between the two effects but reconceptualizes it: rather than reflecting competition between parallel routes, the effects index distinct temporal stages of a serial process, with stem frequency effects arising early and obligatorily and family size effects arising later and conditionally on semantic coherence.

Grainger and colleagues' morpho-orthographic account [CITATION] offers a partially overlapping but architecturally distinct interpretation. On this view, an intermediate level of representation between sublexical letter clusters and whole-word form representations is active in skilled reading, at which morphological constituents are represented as units—not as stored symbolic morphemes, but as learned orthographic patterns that reliably co-occur with semantic content. Stem frequency on this account indexes the strength of representation at this morpho-orthographic level: a high-frequency stem has a more robustly activated intermediate representation, facilitating subsequent access to meaning. Morphological family size, in turn, is reflected in the richness of the semantic representation that a morpho-orthographic unit connects to—a larger family means a more densely structured semantic neighborhood, supporting faster and more confident lexical retrieval. This account differs from the early decomposition account in treating the intermediate representations as learned, gradient, and form-based rather than as the output of a parsing procedure, but like the early decomposition account it assigns the two effects to different levels of a multi-stage architecture.

The most radical reinterpretation is offered by the discriminative learning framework, instantiated in Naive Discriminative Learning [NDL; CITATION] and its successor, Linear Discriminative Learning [LDL; CITATION]. These models share the commitment that morphological processing involves no stored morpheme representations, no decomposition procedure, and no dedicated morphological architecture of any kind. Instead, the system learns a direct mapping from distributional form cues—typically character n-grams—to semantic representations, via error-driven weight adjustment. Within this framework, stem frequency and family size effects are reframed not as diagnostics of separate routes or stages, but as emergent consequences of a single learning mechanism operating over different aspects of distributional structure.

Stem frequency, in NDL and LDL, reflects the accumulated weight strength of the n-gram cues that spell out the base morpheme. Because these cues appear across all surface forms of the base—across every inflected and derived word in which the stem is present—they accumulate training signal proportional to the base's cumulative token frequency. High stem frequency therefore means that the relevant form cues have been more extensively and precisely calibrated, producing stronger and more reliable activation of the associated semantic representation. Critically, this occurs without any discrete stem entry being accessed: the effect is a gradient property of the learned weight matrix rather than the consequence of retrieving a stored morpheme. Morphological family size, by contrast, reflects a different statistical property of the same learning environment: the type diversity of the training signal rather than its token accumulation. A large morphological family exposes the stem's n-gram cues to a wide variety of distinct form-meaning pairs, each sharing the same semantic core but differing in their idiosyncratic features. The error-driven learning dynamics of NDL—and the analogous optimization pressure in LDL's matrix learning—cause the weights for shared cues to converge on the semantic dimensions that are consistent across the family, while suppressing idiosyncratic variation. A large, semantically coherent family therefore produces a more robustly and precisely represented semantic core, facilitating lexical retrieval. The well-documented finding that semantically opaque relatives fail to contribute to the family size effect [CITATION] follows without additional stipulation: an opaque relative has a semantic vector that is distant from the base in distributional space, and therefore provides no consistent reinforcement of the stem's core semantic weights during learning.

The theoretical significance of the NDL/LDL account lies not only in its reframing of each effect individually, but in its claim that both effects arise from a single mechanism. Where dual-route and multi-stage models took the dissociation between stem frequency and family size as evidence for distinct representational levels or processing routes, the discriminative learning framework attributes the dissociation to two aspects of the same distributional statistics—token frequency and type diversity—acting through the same learned weight structure. This is a substantively more parsimonious account, though it comes at the cost of reconceptualizing what the effects are evidence for: rather than revealing the architecture of a morphological processor, they reveal properties of the statistical environment in which the lexicon was learned.

These competing interpretations carry importantly different predictions about the time course of stem frequency and family size effects, about how the system responds to novel morphological structure, and about the role of individual differences in processing efficiency—three dimensions along which the present study is specifically designed to adjudicate.

Individual Differences in Morphological Processing

A complementary perspective on the architecture of morphological processing comes from examining variation across readers. Traditional models of word recognition have typically characterised processing at the level of the skilled reader in general, abstracting over individual differences in favour of a uniform cognitive architecture. Yet if morphological sensitivity reflects the accumulated product of statistical learning rather than the output of a fixed parsing mechanism, readers who differ in the depth or precision of their lexical representations should differ systematically in how they process morphologically complex words—and potentially in which aspects of that processing are most affected.

The Lexical Quality Hypothesis [CITATION] provides a principled account of such differences. On this view, lexical representations vary in quality across readers—that is, in how completely and precisely a word's orthographic, phonological, and semantic constituents are specified and interconnected. High-quality representations enable rapid, automatic activation of a word's sound and meaning from its visual form; low-quality representations are partial or underspecified, producing processing that is slower, less reliable, and more dependent on contextual support. Individual differences in reading and spelling skill can be understood as reflecting differences in lexical quality [CITATION], with more skilled readers possessing more precise orthographic representations that support faster word identification and greater resistance to interference from orthographically similar forms.

Critically, readers do not differ along a single dimension of lexical quality. Andrews and colleagues [CITATION] have demonstrated that individual differences in orthographic and semantic knowledge are dissociable and predict distinct aspects of lexical processing. Readers whose spelling ability exceeds their vocabulary knowledge—an orthographic profile reflecting high precision in sublexical form representations—show robust morphological priming even for semantically opaque pairs such as corner–corn, suggesting that their processing is driven primarily by form-based segmentation at the level of morpho-orthographic structure. Readers whose vocabulary knowledge exceeds their spelling ability—a semantic profile reflecting richer lexical-semantic representations—show morphological priming only for transparent pairs, indicating that their processing is more strongly constrained by meaning. These dissociable profiles suggest that the relative contributions of form-based and meaning-based information to morphological processing are not fixed properties of the recognition system but vary systematically across readers as a function of the type of lexical knowledge they have accumulated.

This dissociation maps directly onto the theoretical contrast between the processing accounts reviewed above. Readers with a strong orthographic profile resemble the system described by morpho-orthographic and decompositional accounts—their processing is driven by sensitivity to sublexical form structure, the primary cue in discriminative learning models and the basis for early automatic decomposition in parsing accounts. Readers with a strong semantic profile resemble the system in which family size and semantic coherence are the dominant influences, consistent with the later semantic stage proposed by multi-stage accounts and with the meaning-space organisation of LDL. Rather than treating this as an either-or contrast between reader types, the present study operationalises these profiles as continuous dimensions derived from a principal components analysis of vocabulary knowledge, spelling ability, and print exposure, yielding two orthogonal components—Language Proficiency and Orthographic Sensitivity—that capture the full range of the dissociation in the sample. If the effects of stem frequency and family size are modulated differently by these two dimensions, and if those modulations emerge at distinct temporal stages in the ERP, this would constitute strong evidence that the theoretical distinction between form-based and meaning-based morphological processing is not merely a property of competing computational models but a genuine axis of variation in the human reading system.

\subsection{Aims and Predictions}

Despite robust evidence that morphological structure shapes visual word recognition, the theoretical debate between structured parsing accounts and discriminative learning frameworks has remained difficult to resolve at the level of group-averaged data. Both classes of model can accommodate the core behavioral effects—facilitation by stem frequency and morphological family size—and both can be reconciled with the broad temporal dissociation between early form-based and later semantic processing. What they cannot equally accommodate are fine-grained differences in when and for whom these effects emerge, and whether the system generalizes morphological regularities productively to novel forms. Our primary aim is therefore to use individual differences as a diagnostic tool: if morphological sensitivity reflects accumulated statistical learning rather than fixed representational primitives, then Orthographic Sensitivity and Language Proficiency should modulate morphological effects in ways that are temporally dissociable—with Orthographic Sensitivity predicting early form-level processing and Language Proficiency predicting later semantic processing—rather than scaling uniformly with general language experience. This predicted dissociation constitutes the study's central theoretical claim, and it is evaluated through two complementary analyses.

The first analysis examines the time course of stem frequency and family size effects in real word recognition. Because the real words in the present study were not matched across morphological variables—reflecting the naturalistic distributional properties of morphologically complex words rather than an orthogonally controlled factorial design—this analysis treats stem frequency, family size, and morphological complexity as continuous predictors in a regression framework for reaction time, examining their independent contributions across the full item set. For ERP data, signal averaging requirements necessitate a different approach: items were median-split into high and low groups separately for stem frequency and family size, and ERPs were averaged within each group to yield stable waveforms for comparison. The two dependent measures therefore carry somewhat different inferential weight: the RT regression speaks to the independent continuous contributions of each variable across the full distributional range, while the ERP analyses speak to the direction and timing of effects at the extremes of each variable's distribution. Multi-stage architectures—whether decompositional [CITATION] or morpho-orthographic [CITATION]—predict that stem frequency and family size should influence processing at distinct temporal stages, with stem frequency modulating earlier components reflecting form-level access and family size exerting its primary influence later when semantic information becomes available. The discriminative learning framework [CITATION], by contrast, treats both effects as arising from a single learned weight structure, predicting no systematic difference in their temporal onset or scalp distribution. We predicted that this temporal dissociation, if present, would be most pronounced in readers high in Language Proficiency—for whom family size should carry particular weight—and that early stem frequency effects would be most strongly modulated by Orthographic Sensitivity.

The second analysis examines generalization to novel morphological structure using morphologically complex nonwords. Because morphologically complex nonwords have no stored whole-word representation, they place maximal pressure on whatever generalization mechanism the system possesses and allow stem-level variables to be examined independently of whole-word familiarity—a control that the real word analysis cannot straightforwardly provide. Morphologically structured nonwords (gasful: real stem + real affix) were compared against matched controls (gasfil: real stem + pseudoaffix), holding stem identity constant while manipulating the morphological legitimacy of the suffix. Rule-governed parsing accounts predict that this manipulation should yield early ERP differences, specifically modulation of the N250, reflecting the detection of constituent structure at the suffix boundary—an effect that should be present categorically wherever a valid parse is licensed and should not vary systematically across readers of comparable literacy. The discriminative learning framework also predicts sensitivity to the suffix manipulation, since real affixes carry stronger and more consistent n-gram cue mappings to morphological outcomes than pseudoaffixes do, but attributes this to the same learned weight structure underlying real word recognition rather than to a dedicated parsing operation. The present design is therefore better positioned to evaluate where in time and for whom suffix legitimacy influences processing than to adjudicate between categorical and graded generalization directly. Critically, if Orthographic Sensitivity modulates the magnitude or onset of the N250 response to the suffix contrast, this would implicate form-cue learning rather than a fixed parsing operation, since the output of a rule-governed parser should be invariant across readers who have all clearly acquired basic morphological knowledge.

Taken together, the two analyses ask whether the distinction between form-based and meaning-based morphological processing—long debated as a property of competing computational architectures—is also a genuine axis of variation in the human reading system. If Orthographic Sensitivity selectively modulates early ERP effects of stem frequency and suffix structure while Language Proficiency selectively modulates later N400 effects of family size, the data would support a view in which different types of accumulated lexical knowledge contribute to qualitatively distinct stages of morphological recognition, with consequences for how individual differences in reading should be theorised and measured.

